\section{Transfer bytes and establish EOF before parsing}
The loader first checks that its unsigned capacity lies from one through 266240.
Then open the input, accumulate actual positive read counts and continue to EOF.
When the caller buffer becomes exactly full, use a private one-byte probe to distinguish
an exact fit from more input. The probe's byte must not be stored beyond the caller
buffer. Close the descriptor before parsing the accepted complete bytes and publishing
metadata. A failed load may retain its accepted buffer prefix; the descriptor stays
unchanged.

A read error, including an error at the probe, returns I/O status three. An extra byte
past capacity returns two. Once a complete file has been established, an unsupported
or malformed container returns one. For example, a valid forty-four-byte empty-audio
file with capacity 43 fails the loader's bound, even though its representation would
pass parsing in a large enough buffer. Do not replace that earlier observation with
a later format guess.

The saver validates the selected complete stream before opening output in truncate
mode. Advance by actual positive short writes, reject zero progress or a negative
result, and close every acquired descriptor once on each return path. A later failure
cannot undo truncation or an accepted prefix. The current educational write syscall
also reads one byte beyond its requested interval, including on short or error returns.
Provide a readable guard outside the file extent, but never add it to the file count.
